\section{Related Work}
\label{sec:related_work}

\subsection{Low-light image enhancement}

Low-light image enhancement has been studied through both physical image formation models and data-driven restoration networks. Retinex theory models an image as the interaction of reflectance and illumination~\cite{land1971retinex}, motivating decomposition-based methods that estimate and manipulate illumination maps. RetinexNet~\cite{wei2018deepretinex} and KinD~\cite{zhang2019kind} extend this idea with neural decompositions and enhancement modules, while Zero-DCE~\cite{guo2020zerodce} and Zero-DCE++~\cite{li2021zerodcepp} formulate enhancement as curve estimation without paired supervision. More recent transformer-based methods, including Retinexformer~\cite{cai2023retinexformer}, substantially improve RGB restoration fidelity. A 2024--2026 refresh shows that the field is also moving toward state-space and exposure-correction architectures~\cite{weng2024mamballie,dong2024ecmamba,li2024osmamba}, unfolding and Retinex-consistency designs~\cite{he2025unfoldir,xu2025consistretinex}, and diffusion or flow-based restoration~\cite{xu2025boostingfidelityfo,wei2025illumflow}. These methods establish strong enhancement baselines, and our experiments confirm that Retinexformer remains far ahead of our RLEF variants on LOL-v2 synthetic PSNR.

Our work is not positioned as a stronger restoration network. Instead, it studies a failure mode that standard restoration metrics can hide: a method may improve RGB fidelity while learning an exposure field whose spatial structure is physically inconsistent. This shifts the question from ``which enhancer has the highest PSNR?'' to ``when an enhancer claims to learn an exposure field, is that field identifiable and correctly evaluated?''

\subsection{Exposure fields and physical interpretability}

Many physically motivated enhancement methods use illumination maps, exposure maps, transmission-like fields, or Retinex factors as interpretable internal variables. Such variables are attractive because they promise more than black-box restoration: they suggest that the model has discovered a meaningful correction structure. Recent work has also made non-uniform exposure explicit, including dual uniform/non-uniform LLIE~\cite{perez-zarate2024alen}, spatially non-uniform exposure correction~\cite{li2026rethinkingexposure}, log-domain intensity--chromaticity decoupling~\cite{bai2026rethinkinglowlight}, and illumination-field or Gaussian light-field modeling~\cite{chen2026aigsnet,chen2026gaussianlightfield}. These trends make identifiability more important, not less. In gradient-domain or Poisson-like objectives, adding a constant to an exposure field preserves local gradients, so absolute gauge and spatial shape can be confounded. A field may therefore appear plausible under one metric while being wrong under another.

GIR-Field addresses this issue by separating the field into a gauge-free shape $S$ and an image-level gauge $\mu$. This is a modest but important shift: rather than treating exposure error as a single scalar, we ask whether an error comes from global gauge shift, local shape distortion, or output restoration failure. This decomposition is the basis for both our training objective and our evaluation protocol.

\subsection{Benchmarks and metric design for enhancement}

Paired datasets such as LOL-v2 and low-light raw-image benchmarks~\cite{chen2018sid} support full-reference evaluation with PSNR and SSIM, while unpaired real collections are often evaluated with no-reference metrics or visual inspection. Recent challenge and evaluation work reinforces that LLIE assessment is becoming broader but remains mostly output-centered: NTIRE 2025/2026 challenge reports emphasize restoration efficiency and quality~\cite{liu2025ntire2025challenge,yan2026ntire2026challenge}, multi-intensity evaluation probes robustness across illumination levels~\cite{pilligua2025evaluatinglowlight}, and low-light quality assessment studies artifacts such as noise, color shift, structural damage, and over-exposure~\cite{sun2026leiqassessor}. These metrics are useful but incomplete for exposure-field models. PSNR can reward a model that reconstructs RGB values well while giving no guarantee that the internal exposure map is physically meaningful. No-reference scores can provide additional evidence, but they are not stable enough to support a main superiority claim.

UEFB-G is designed for this gap. It does not replace PSNR; instead, it prevents PSNR from being the only ranking signal. Output-only metrics remain available for any method, including Retinexformer, Zero-DCE++, and KinD++. Internal-field metrics are reported only when a method explicitly predicts an exposure field. This distinction avoids unfairly penalizing black-box baselines for not exposing internal variables while still testing the physical claims of exposure-field methods.

\subsection{Calibration, selective prediction, and recoverability risk}

The recoverability question is related to selective prediction and calibration: an enhancer should know where enhancement is beneficial, harmful, or uncertain. Prior work in calibrated prediction emphasizes not only the accuracy of a score but also whether the score's confidence matches empirical risk. In our setting, DGB-style routing experiments suggested that route usage alone is not sufficient: sending more samples through a restoration or safe path does not guarantee better output quality.

We therefore treat recoverability as a diagnostic audit rather than a solved module. Our lightweight risk probe reaches AUC around 0.766, showing that harmful enhancement is predictable from image and field statistics, but ECE around 0.281 shows that calibration remains weak. This result supports future work on recoverability-aware enhancement, but it does not justify a main method claim.

\subsection{Positioning of this work}

Compared with restoration-centered LLIE methods, GIR-Field contributes a claim-calibrated evaluation and mechanism analysis. Compared with Retinex-style decomposition, it focuses on the gauge ambiguity of the exposure field and explicitly separates global gauge from spatial shape. Compared with benchmark-only work, it provides a concrete mechanism---centered E-shape calibration---that improves gauge-free shape correlation and paired fidelity in a formal frozen-evidence audit. The resulting paper is therefore best read as a mechanism-and-benchmark contribution rather than a SOTA restoration paper.
